\chapter{Pencacahan}\label{ch:b1:counting}

Mencacah \omterm{def:b1:counting:card}{himpunan hingga} terdengar sederhana --- dan dengan cepat
menjadi halus. Bab ini mendefinisikan \omterm{def:b1:counting:card}{kardinalitas} secara benar (lewat
bijeksi, sejalan dengan \cref{ch:b1:logic}), menegakkan segelintir
asas pencacahan yang darinya segala sesuatu mengikut, lalu menurunkan
hasil cacah yang klasik: daftar, \omterm{def:b1:counting:objects}{permutasi}, \omterm{def:b1:logic:sets}{himpunan} bagian, \omterm{def:b1:counting:objects}{koefisien binomial}.

\section{Kardinalitas himpunan hingga}

\begin{definition}[Himpunan hingga, kardinalitas]\label{def:b1:counting:card}
Untuk $n \in \N^*$, tulis $\intint{1}{n} = \{1, 2, \dots,
n\}$. \omterm{def:b1:logic:sets}{Himpunan} $E$ disebut \emph{hingga}\index{himpunan hingga} bila $E = \emptyset$
atau ada bijeksi dari $\intint{1}{n}$ pada $E$ untuk
suatu $n \in \N^*$; bilangan $n$ ini tunggal (\cref{thm:b1:counting:welldef})
dan merupakan \emph{kardinalitas}\index{kardinalitas} dari $E$, ditulis
$\abs{E}$ (dengan $\abs{\emptyset} = 0$).
\end{definition}

\begin{theorem}[Kardinalitas terdefinisi dengan baik]\label{thm:b1:counting:welldef}
Jika $m \neq n$, tidak ada bijeksi dari $\intint{1}{m}$
pada $\intint{1}{n}$. Lebih tepatnya, jika $m > n$ maka tidak ada
injeksi dari $\intint{1}{m}$ ke dalam $\intint{1}{n}$.
\end{theorem}

\begin{proof}
Kita buktikan dengan induksi pada $n$ \omterm{def:b1:logic:statement}{pernyataan} berikut: \emph{untuk setiap $m > n$,
tidak ada injeksi $\intint{1}{m} \to \intint{1}{n}$}. Untuk $n = 0$ sasarannya kosong sedangkan $m \geq 1$: tidak ada \omterm{def:b1:logic:map}{pemetaan}
sama sekali. Andaikan \omterm{def:b1:logic:statement}{pernyataan} itu berlaku untuk $n$, dan andaikan $f \colon
\intint{1}{m} \to \intint{1}{n+1}$ sebuah
injeksi dengan $m > n + 1$. Bila nilai $n + 1$ tidak tercapai, $f$
adalah injeksi ke dalam $\intint{1}{n}$, bertentangan dengan
hipotesis induksinya. Bila tercapai, $f(a) = n + 1$ untuk tepat satu $a$;
tukarkan $f(a)$ dengan $f(m)$ (secara formal: komposisikan dengan transposisi
kedua nilai itu), sehingga injeksi baru $g$ memenuhi $g(m) = n + 1$.
Maka pembatasan $g$ pada $\intint{1}{m-1}$ adalah
injeksi ke dalam $\intint{1}{n}$ dengan $m - 1 > n$ ---
kontradiksi lagi.
\end{proof}

\begin{corollary}[Prinsip sarang merpati]\label{cor:b1:counting:pigeonhole}
Jika $\abs{E} > \abs{F}$, tak ada \omterm{def:b1:logic:map}{pemetaan} $f \colon E \to F$ yang \omterm{def:b1:logic:inj}{injektif}:
ada dua unsur $E$ yang berbagi peta\index{prinsip sarang merpati}.
\end{corollary}

\begin{proof}
Tulis $\abs E = m$, $\abs F = n$ dengan $m > n$, lalu pilih bijeksi
$u \colon \intint1m \to E$ dan $v \colon F \to \intint1n$. Jika $f$
\omterm{def:b1:logic:inj}{injektif}, maka $v \circ f \circ u$ akan menjadi injeksi dari
$\intint1m$ ke dalam $\intint1n$ (komposisi injeksi,
\cref{prop:b1:logic:comp}), bertentangan dengan
\cref{thm:b1:counting:welldef}.
\end{proof}

\begin{remark}[Selingan: mengapa ada penukaran pada bukti teorema itu?]\label{rem:b1:counting:swapinterlude}
Bukti \cref{thm:b1:counting:welldef} memuat langkah cerdik pertama
bab ini, yang layak diputar ulang perlahan.
Rintangannya: untuk memakai hipotesis induksi kita ingin
menghapus titik terakhir $m$ pada sumbernya \emph{dan} titik terakhir
$n+1$ pada sasarannya, tetapi $f$ mungkin mengirim titik lain $a$ ke $n
+ 1$, dan menghapus titik sasaran itu lalu merusak \omterm{def:b1:logic:map}{pemetaan}
di tempat lain. Obatnya: komposisikan $f$ dengan transposisi kedua
\emph{nilai} $f(a)$ dan $f(m)$ --- sebuah bijeksi pada
sasarannya, jadi keinjektifan terjaga --- setelah itu
nilai $n + 1$ yang merepotkan tadi duduk pada posisi $m$ yang tak berbahaya, dan
kedua penghapusan menjadi bersih. Pola ``normalkan dulu, baru potong''
ini berulang: begitulah rekursi \omterm{def:b1:counting:objects}{permutasi} kacau mengalihkan
$\sigma^{-1}(n+1)$ pada soal akhir pekan bab ini, dan
begitulah \omterm{def:b1:counting:objects}{permutasi} ditambal di sepanjang
soal \cref{ch:b1:structures} tentang grup simetri.
\end{remark}

\begin{proposition}[Injeksi, surjeksi dan kardinalitas]\label{prop:b1:counting:injsur}
Misalkan $E, F$ \omterm{def:b1:counting:card}{himpunan hingga} dengan $\abs{E} = \abs{F}$, dan $f \colon E
\to F$. Maka
\[
f \text{ injektif} \iff f \text{ surjektif} \iff f \text{ bijektif}.
\]
\end{proposition}

\begin{proof}
Andaikan $f$ \omterm{def:b1:logic:inj}{injektif}. Maka $f$ adalah bijeksi dari $E$ pada $f(E)$,
sehingga $\abs{f(E)} = \abs{E} = \abs{F}$. Jika $f(E)$ melewatkan sebuah titik $y_0$ pada
$F$, maka $f$ akan menjadi injeksi dari $E$ ke dalam $F \setminus \{y_0\}$,
yaitu \omterm{def:b1:logic:sets}{himpunan} berkardinalitas $\abs{F} - 1 < \abs{E}$ --- mustahil menurut
prinsip sarang merpati. Jadi $f(E) = F$: $f$ \omterm{def:b1:logic:inj}{surjektif}, sehingga
\omterm{def:b1:logic:inj}{bijektif}.

Andaikan $f$ \omterm{def:b1:logic:inj}{surjektif}. Pilih untuk setiap $y \in F$ satu \omterm{def:b1:logic:map}{prapeta} $s(y)
\in E$; maka $f \circ s = \mathrm{id}_F$, jadi $s$ \omterm{def:b1:logic:inj}{injektif}
(\cref{prop:b1:logic:comp}). Menurut paragraf sebelumnya yang diterapkan pada $s$
(\omterm{def:b1:counting:card}{kardinalitas} keduanya sama), $s$ \omterm{def:b1:logic:inj}{bijektif}. Dari $f \circ s =
\mathrm{id}_F$ kita peroleh $f = \mathrm{id}_F \circ s^{-1} = s^{-1}$, jadi
$f$ \omterm{def:b1:logic:inj}{bijektif}. Akhirnya, \omterm{def:b1:logic:map}{pemetaan} \omterm{def:b1:logic:inj}{bijektif} menurut definisinya sekaligus
\omterm{def:b1:logic:inj}{injektif} dan \omterm{def:b1:logic:inj}{surjektif}, dan itu menutup daur implikasinya.
\end{proof}

\begin{example}[Kehinggaan itu penting]\label{ex:b1:counting:finiteonly}
Pada \omterm{def:b1:logic:sets}{himpunan} yang \emph{\omterm{def:b1:counting:card}{hingga}}, \cref{prop:b1:counting:injsur} adalah
jalan pintas yang ampuh: setiap \omterm{def:b1:logic:map}{pemetaan} \omterm{def:b1:logic:inj}{injektif} dari $E$ ke dirinya sendiri
otomatis menjadi \omterm{def:b1:counting:objects}{permutasi} $E$ --- separuh dari kebijektifan datang cuma-cuma.
Kedua implikasi itu runtuh pada \omterm{def:b1:logic:sets}{himpunan} tak \omterm{def:b1:counting:card}{hingga}: $n \mapsto n + 1$
\omterm{def:b1:logic:inj}{injektif} dari $\N$ ke $\N$ tetapi melewatkan $0$, dan \omterm{def:b1:logic:map}{pemetaan} $\N \to \N$
yang mengirim $0 \mapsto 0$ dan $n \mapsto n - 1$ untuk $n \geq 1$
\omterm{def:b1:logic:inj}{surjektif} tetapi tidak \omterm{def:b1:logic:inj}{injektif}. Setiap kali proposisi ini dipanggil,
hipotesis kehinggaannya sedang bekerja sungguhan --- tema yang ditelusuri
soal akhir pekan \cref{ch:b1:logic} dari sisi yang berlawanan,
tempat \omterm{def:b1:logic:sets}{himpunan} tak \omterm{def:b1:counting:card}{hingga} persis adalah \omterm{def:b1:logic:sets}{himpunan} yang punya pemetaan-diri semacam itu.
\end{example}

\begin{example}[Separuh pekerjaan, cuma-cuma]\label{ex:b1:counting:mod7}
Tinjau \omterm{def:b1:logic:map}{pemetaan} $f$ pada $\{0, 1, \dots, 6\}$ yang mengirim $k$ ke
sisa pembagian $3k$ oleh $7$; tabel nilainya adalah
\[
0,\ 3,\ 6,\ 2,\ 5,\ 1,\ 4 .
\]
Apakah $f$ sebuah bijeksi? Keinjektifan saja sudah cukup
(\cref{prop:b1:counting:injsur}): jika $3k$ dan $3k'$ bersisa sama,
maka $7$ membagi $3(k - k')$, dan karena $7$ prima serta
tidak membagi $3$, ia membagi $k - k'$ (lema Euclid, yang di sini dipakai
pada taraf sekolah menengah dan dibuktikan pada \cref{ch:b1:arith}); dengan
$\abs{k - k'} \leq 6$ hal ini memaksa $k = k'$. Kesurjektifan datang
cuma-cuma --- tak perlu menyelesaikan $3k \equiv c$ untuk setiap $c$, meskipun
tabelnya membenarkan bahwa setiap nilai muncul tepat sekali. Jalan pintas itu
adalah kuda beban: ia membuktikan keterbalikan perkalian
modular (\cref{ch:b1:arith}), menggerakkan pemasangan pada
teorema Wilson, dan kembali dalam aljabar linear sebagai ``endomorfisma
ruang berdimensi \omterm{def:b1:counting:card}{hingga} bersifat \omterm{def:b1:logic:inj}{injektif} jika dan hanya jika
\omterm{def:b1:logic:inj}{surjektif}'' (\cref{ch:b1:findim}).
\end{example}

\section{Asas pencacahan}

\begin{proposition}[Kaidah jumlah dan kaidah hasil kali]\label{prop:b1:counting:rules}
Misalkan $E, F$ \omterm{def:b1:counting:card}{himpunan hingga}.
\begin{enumerate}
  \item Jika $E \cap F = \emptyset$, maka $\abs{E \cup F} = \abs{E} +
        \abs{F}$; lebih umum, untuk \omterm{thm:b1:logic:partition}{partisi} $E$ atas potongan
        $E_1, \dots, E_k$ berlaku $\abs{E} = \sum_i \abs{E_i}$.
  \item Secara umum, $\abs{E \cup F} = \abs{E} + \abs{F} - \abs{E \cap
        F}$.
  \item $\abs{E \times F} = \abs{E} \times \abs{F}$.
  \item \omterm{def:b1:logic:sets}{Himpunan} $F^E$ berisi semua \omterm{def:b1:logic:map}{pemetaan} dari $E$ ke $F$ memenuhi
        $\abs{F^E} = \abs{F}^{\abs{E}}$.
  \item $\abs{\mathcal{P}(E)} = 2^{\abs{E}}$.
\end{enumerate}
\end{proposition}

\begin{proof}
(1) Sambungkan pencacahannya: jika $E = \{x_1, \dots, x_m\}$ dan $F =
\{y_1, \dots, y_n\}$ tanpa pengulangan, maka $x_1, \dots, x_m, y_1,
\dots, y_n$ mencacah $E \cup F$ tanpa pengulangan (karena saling lepas).
Induksi memperluasnya ke $k$ potongan.

(2) $E \cup F$ adalah gabungan saling lepas dari $E$ dan $F \setminus E$, dan
$F$ adalah gabungan saling lepas dari $F \cap E$ dan $F \setminus E$; jadi
$\abs{E \cup F} = \abs{E} + \abs{F \setminus E} = \abs{E} + \abs{F} -
\abs{E \cap F}$.

(3) $E \times F$ adalah gabungan saling lepas, atas $x \in E$, dari \omterm{def:b1:logic:sets}{himpunan}
$\{x\} \times F$, yang masing-masing berkardinalitas $\abs{F}$; terapkan (1).

(4) Sebuah \omterm{def:b1:logic:map}{pemetaan} dari $E = \{x_1, \dots, x_m\}$ ke $F$ persis sama dengan pemilihan
tupel-$m$ $(f(x_1), \dots, f(x_m)) \in F^m$; padanan
ini adalah bijeksi, dan $\abs{F^m} = \abs{F}^m$ menurut (3) beserta
induksi.

(5) \omterm{def:b1:logic:sets}{Himpunan} bagian $E$ berpadanan secara \omterm{def:b1:logic:inj}{bijektif} dengan \omterm{def:b1:logic:map}{pemetaan} $E \to \{0, 1\}$
(kirim $A$ ke fungsi indikatornya); terapkan (4).
\end{proof}

\begin{example}[Mencacah lewat komplemen]\label{ex:b1:counting:complement}
Berapa banyak kode PIN $4$ angka (angka $0$--$9$, urutan penting,
pengulangan dibolehkan) yang memuat \emph{sekurang-kurangnya satu} angka berulang?
Mencacahnya secara langsung berarti kita harus menjungkirkan kasus ``tepat satu pasang,
dua pasang, kembar tiga, kembar empat'' --- lima konfigurasi yang saling
bertindihan. Cacah komplemennya saja: seluruh kode berjumlah
$10^4 = 10\,000$ (kaidah hasil kali), sedangkan kode dengan empat angka berbeda
berjumlah $10 \times 9 \times 8 \times 7 = 5\,040$
(susunan-$4$), jadi jawabannya adalah
\[
10^4 - 10 \cdot 9 \cdot 8 \cdot 7 = 10\,000 - 5\,040 = 4\,960 .
\]
Hampir separuh dari semua PIN mengulang sebuah angka. Inti gagasannya:
setiap kali sebuah cacahan dirumuskan dengan ``sekurang-kurangnya'' atau
``tidak semua'', cobalah komplemennya lebih dulu --- kaidah jumlah menjamin
$\abs{A} = \abs{E} - \abs{\overline A}$, dan komplemennya
sering berupa satu konfigurasi yang bersih.
\end{example}

\begin{example}[Lintasan pada kisi]\label{ex:b1:counting:paths}
Cacah lintasan terpendek dari sudut $(0,0)$ ke sudut
$(4, 3)$ sebuah kisi, dengan hanya boleh melangkah satu petak ke kanan (K)
atau satu petak ke atas (A) setiap kali. Setiap lintasan semacam itu
menempuh tepat $7$ langkah, yang
$4$ di antaranya K dan $3$ di antaranya A; sebaliknya, setiap kata sepanjang $7$ atas
huruf K, A dengan empat huruf K menggambarkan tepat satu lintasan. Jadi lintasan
berpadanan secara \omterm{def:b1:logic:inj}{bijektif} dengan pilihan posisi
huruf K:
\[
\binom{7}{4} = 35 .
\]
Inti gagasannya ada pada \emph{penyandiannya}: cacahan itu menjadi sepele
begitu setiap lintasan diterjemahkan menjadi sebuah kata, yakni sebuah
\omterm{def:b1:logic:sets}{himpunan} bagian posisi --- satu lagi contoh dari semboyan bahwa cacahan
yang benar adalah bijeksi yang menyamar (\cref{met:b1:counting:which}).
\end{example}

\begin{omfigure}
\begin{tikzpicture}[scale=0.85]
  \draw[gray!60] (0,0) grid (4,3);
  \draw[->, thin] (-0.4,0) -- (4.6,0);
  \draw[->, thin] (0,-0.4) -- (0,3.6);
  \draw[omDef, very thick]
    (0,0) -- (1,0) -- (1,1) -- (3,1) -- (3,2) -- (4,2) -- (4,3);
  \fill (0,0) circle (0.07) node[below left] {$(0,0)$};
  \fill (4,3) circle (0.07) node[above right] {$(4,3)$};
  \node[omDef, below] at (0.5,0) {\small K};
  \node[omDef, left] at (1,0.5) {\small A};
  \node[omDef, below] at (1.5,1) {\small K};
  \node[omDef, below] at (2.5,1) {\small K};
  \node[omDef, left] at (3,1.5) {\small A};
  \node[omDef, below] at (3.5,2) {\small K};
  \node[omDef, left] at (4,2.5) {\small A};
\end{tikzpicture}

{\small Salah satu dari $\binom74 = 35$ lintasan terpendek dari $(0,0)$ ke
$(4,3)$: lintasan yang digambar menyandikan kata KAKKAKA, yakni pilihan
posisi $\{1,3,4,6\}$ bagi huruf K di antara ketujuh
langkahnya.}
\end{omfigure}

\section{Daftar, permutasi, himpunan bagian}

\begin{definition}[Susunan, permutasi, kombinasi]\label{def:b1:counting:objects}
Misalkan $E$ \omterm{def:b1:logic:sets}{himpunan} dengan $\abs{E} = n$ dan $0 \leq k \leq n$.
\begin{itemize}
  \item \emph{Susunan-$k$} dari $E$ adalah tupel-$k$ \omterm{def:b1:logic:inj}{injektif} berisi
        unsur $E$ (pemilihan terurut tanpa pengulangan);
  \item \emph{permutasi}\index{permutasi} dari $E$ adalah bijeksi
        dari $E$ ke dirinya sendiri --- setara dengan susunan-$n$;
  \item \emph{kombinasi-$k$} adalah \omterm{def:b1:logic:sets}{himpunan} bagian $E$ beranggota $k$ unsur
        (pemilihan tak terurut tanpa pengulangan). Banyaknya
        ditulis $\binom{n}{k}$, dibaca ``$n$ pilih $k$''
        \index{koefisien binomial}.
\end{itemize}
\end{definition}

\begin{theorem}[Ketiga cacahan itu]\label{thm:b1:counting:counts}
Dengan $n = \abs{E}$ dan $0 \leq k \leq n$:
\begin{enumerate}
  \item banyaknya susunan-$k$ dari $E$ adalah $n (n-1) \cdots
        (n-k+1) = \dfrac{n!}{(n-k)!}$;
  \item banyaknya \omterm{def:b1:counting:objects}{permutasi} $E$ adalah $n!$;
  \item $\dbinom{n}{k} = \dfrac{n!}{k!\,(n-k)!}$.
\end{enumerate}
\end{theorem}

\begin{proof}
(1) Pilih koordinat pertama ($n$ cara), lalu yang kedua ($n - 1$
pilihan tersisa), \dots, lalu yang ke-$k$ ($n - k + 1$ pilihan).
Secara formal, berinduksilah pada $k$. Untuk $k = 1$ ada $n$ tupel
\omterm{def:b1:logic:inj}{injektif} bersuku satu. Andaikan cacahannya berlaku untuk $k - 1$. Setiap
susunan-$k$ $(x_1, \dots, x_k)$ diperoleh dari tepat satu
susunan-$(k-1)$ --- yaitu pemenggalannya $(x_1, \dots, x_{k-1})$
--- dengan menambahkan koordinat terakhir di luar $\{x_1, \dots,
x_{k-1}\}$, yang untuknya tepat $n - (k - 1)$ nilai tersedia.
Jadi susunan-$k$ terpartisi, lewat pemenggalan, atas
kelas-kelas berukuran sama $n - k + 1$ yang terindeks oleh
susunan-$(k-1)$, dan kaidah jumlah memberikan
\[
\frac{n!}{(n-k+1)!}\;(n - k + 1) = \frac{n!}{(n-k)!} .
\]

(2) adalah (1) dengan $k = n$.

(3) Setiap \omterm{def:b1:logic:sets}{himpunan} bagian beranggota $k$ terurutkan menjadi $k!$ susunan-$k$ yang berbeda, dan
setiap susunan-$k$ muncul dari tepat satu \omterm{def:b1:logic:sets}{himpunan} bagian: jadi
$\frac{n!}{(n-k)!} = \binom nk \cdot k!$.
\end{proof}

\begin{example}[Meja bundar: membagi habis simetrinya]\label{ex:b1:counting:roundtable}
Dengan berapa cara $n$ tamu dapat duduk mengelilingi meja bundar, bila dua
susunan dianggap sama ketika setiap tamu mempunyai tetangga kiri dan
kanan yang sama --- yakni sampai rotasi? Setiap susunan melingkar
berpadanan dengan tepat $n$ susunan lurus (potong lingkarannya di salah satu
dari $n$ tempat), jadi $n!$ urutan lurus itu melebur dalam kelompok
beranggota $n$:
\[
\frac{n!}{n} = (n-1)! \quad\text{susunan melingkar.}
\]
Setara dengan itu: dudukkan satu tamu istimewa di mana saja (mematikan
kebebasan rotasinya), lalu urutkan $n - 1$ tamu sisanya
searah jarum jam. Untuk $n = 6$: ada $120$ meja. Kedua penyelesaian itu
memperlihatkan dua obat baku bagi pencacahan berlebih: bagilah dengan
banyaknya pengulangan yang persis, atau \emph{patahkan simetrinya} dengan
memaku satu objek. Keduanya menuntut ukuran kelompok pengulangan itu sama
bagi setiap konfigurasi --- syarat yang juga dipakai bukti rumus
$\binom nk = \frac{n!}{k!\,(n-k)!}$ di atas, dengan
$k!$ menggantikan $n$.
\end{example}

\begin{example}[Menambahkan satu kendala]\label{ex:b1:counting:constraint}
Melanjutkan meja bundar tadi: di antara $(n-1)!$ meja berisi $n \geq
3$ tamu, berapa banyak yang mendudukkan dua tamu tertentu $A$ dan $B$ \emph{terpisah}
(tidak berdampingan)? Cacah komplemennya. Meja yang $A$ dan $B$-nya
duduk bersama: rekatkan keduanya menjadi satu blok --- $n - 1$ objek
mengelilingi meja, yakni $(n-2)!$ susunan melingkar --- lalu
urutkan pasangan itu di dalam bloknya ($2$ cara): ada $2\,(n-2)!$ meja
yang berdampingan. Jadi
\[
(n-1)! - 2\,(n-2)! = (n-2)!\,\bigl((n - 1) - 2\bigr)
= (n-3)\,(n-2)!
\]
meja membuat keduanya terpisah. Pemeriksaan kewajaran: $n = 3$ memberikan $0$ (mengelilingi
segitiga, semua orang bersentuhan) dan $n = 4$ memberikan $2$,
yang mudah didaftar dengan tangan. Kiat perekatan --- perlakukan blok yang
dipaksakan sebagai satu objek, lalu cacah susunan di dalamnya --- adalah
obat baku bagi kendala kebersebelahan, baik lurus maupun melingkar.
\end{example}

\begin{proposition}[Kesamaan dasar]\label{prop:b1:counting:identities}
Untuk $0 \leq k \leq n$:
\[
\binom{n}{k} = \binom{n}{n-k},
\qquad
\binom{n}{k} = \binom{n-1}{k-1} + \binom{n-1}{k}
\quad (1 \leq k \leq n-1),
\qquad
\sum_{k=0}^{n} \binom{n}{k} = 2^n .
\]
\end{proposition}

\begin{proof}
Kesamaan pertama: $A \mapsto E \setminus A$ adalah bijeksi antara
\omterm{def:b1:logic:sets}{himpunan} bagian beranggota $k$ dan \omterm{def:b1:logic:sets}{himpunan} bagian beranggota $(n-k)$. Kaidah Pascal\index{kaidah Pascal}:
tetapkan sebuah unsur $a \in E$; \omterm{def:b1:logic:sets}{himpunan} bagian beranggota $k$ terbelah atas yang memuat
$a$ (pilih $k - 1$ unsur lainnya: $\binom{n-1}{k-1}$) dan yang menghindari
$a$ ($\binom{n-1}{k}$). Kesamaan ketiga: kedua ruasnya mencacah semua \omterm{def:b1:logic:sets}{himpunan} bagian
$E$, dan di ruas kiri cacahan itu dipilah menurut ukurannya
(\cref{prop:b1:counting:rules}~(1) dan (5)).
\end{proof}

\begin{theorem}[Teorema binomial]\label{thm:b1:counting:binomial}
Untuk setiap $a, b$ dalam sebuah ring komutatif (misalnya $\R$ atau $\C$) dan $n \in \N$:
\[
(a + b)^n = \sum_{k=0}^{n} \binom{n}{k} a^k b^{\,n-k} .
\]
\end{theorem}

\begin{proof}
Menjabarkan $(a+b)(a+b)\cdots(a+b)$ secara distributif menghasilkan satu suku bagi setiap
pilihan, pada masing-masing faktor, antara $a$ atau $b$: suku $a^k b^{n-k}$
muncul sekali untuk setiap cara memilih $k$ faktor mana di antara $n$ faktor itu yang
menyumbang $a$ --- yaitu sebanyak $\binom nk$ kali. (Cara lain: berinduksi
pada $n$ memakai kaidah Pascal.)
\end{proof}

\begin{example}\label{ex:b1:counting:binomial}
Dua pengkhususan klasik: $a = b = 1$ memulihkan $\sum_k \binom nk
= 2^n$; $a = -1$, $b = 1$ memberikan $\sum_{k} (-1)^k \binom nk = 0$ untuk
$n \geq 1$: di antara \omterm{def:b1:logic:sets}{himpunan} bagian sebuah \omterm{def:b1:logic:sets}{himpunan} tak kosong, tepat separuhnya
berkardinalitas genap.
\end{example}

\begin{example}[Satu kesamaan, dua bukti]\label{ex:b1:counting:threen}
Pengkhususan $a = 2$, $b = 1$ pada teorema binomial berbunyi
\[
\sum_{k=0}^{n} \binom nk\,2^k = 3^n .
\]
Berikut kesamaan yang sama tanpa aljabar sama sekali. Ruas kanan
mencacah kata sepanjang $n$ atas abjad $\{0, 1, 2\}$
(kaidah hasil kali). Golongkan setiap kata menurut \omterm{def:b1:logic:sets}{himpunan} $K$ berisi posisi
yang memuat huruf tak nol: memilih $K$ dengan $\abs K = k$ menelan biaya
$\binom nk$, lalu masing-masing posisi $K$ secara bebas memuat $1$
atau $2$: ada $2^k$ cara. Kaidah jumlah atas $k$ memberikan ruas kiri.
Selain kenikmatan melihat keduanya cocok, kedua bukti itu punya keutamaan
yang berbeda: yang aljabar dapat diperluas ke sebarang nilai $a$, sedangkan yang
kombinatorial \emph{menjelaskan} rumusnya dan menyesuaikan diri pada
kendala (larang huruf $2$ pada posisi terakhir, misalnya) yang
tak tertangkap oleh substitusi apa pun. Menjaga kedua teknik itu tetap hidup adalah
keterampilan praktis yang dilatih bab ini.
\end{example}

\begin{method}[Cacahan mana yang berlaku?]\label{met:b1:counting:which}
Sebelum menghitung, jawablah dua pertanyaan tentang pemilihannya: apakah
\emph{urutan} penting, dan apakah \emph{pengulangan} dibolehkan?
\begin{center}
\renewcommand{\arraystretch}{1.4}
\begin{tabular}{l|c|c}
 & urutan penting & urutan tidak penting \\
\hline
tanpa pengulangan & $\dfrac{n!}{(n-k)!}$ & $\dbinom{n}{k}$ \\[6pt]
pengulangan dibolehkan & $n^k$ & (\cref{exo:b1:counting:10}) \\
\end{tabular}
\end{center}
Lalu carilah bijeksi atau \omterm{thm:b1:logic:partition}{partisi} yang menyusutkan masalahnya menjadi
cacahan model itu; cacahan yang benar adalah bijeksi yang menyamar.
\end{method}

\begin{remark}[Jebakan yang lazim dalam pencacahan]\label{rem:b1:counting:pitfalls}
\begin{enumerate}
  \item \emph{Menjumlahkan kasus yang tidak saling lepas.} Kaidah jumlah
        menuntut sebuah \omterm{thm:b1:logic:partition}{partisi}; jika sebuah konfigurasi dapat memenuhi
        dua kasus sekaligus, ia tercacah dua kali --- obatnya adalah
        inklusi--eksklusi (\cref{thm:b1:counting:inclexcl}) atau
        pemilahan kasus yang lebih halus.
  \item \emph{Terurut lawan tak terurut.} Memilih ``panitia beranggota
        dua'' adalah $\binom n2$, bukan $n(n-1)$: putuskan \emph{sebelum
        menghitung} apakah pemilihannya membawa urutan, dan bila
        cacahan terurut lebih mudah, bagilah dengan banyaknya
        pengurutan pada akhirnya --- tetapi hanya bila setiap objek tak
        terurut muncul dari \emph{sama} banyak objek terurut.
  \item \emph{Pilihan bertahap yang tidak saling bebas.}
        Kaidah hasil kali menuntut banyaknya pilihan pada setiap tahap
        tidak bergantung pada pilihan sebelumnya. ``Pilih seorang
        kapten, lalu seorang wakil kapten yang berbeda'' aman ($n(n-1)$);
        ``pilih dua pemain yang cocok satu sama lain'' sama sekali bukan
        hasil kali dua tahap.
  \item \emph{Pencacahan ganda karena konstruksinya.} Membangun setiap
        objek dua kali --- misalnya mencacah tangan berisi
        \emph{sekurang-kurangnya} satu raja sebagai (pilih satu raja)
        $\times$ (pilih $4$ kartu lagi) --- melebihkan cacahan tangan
        yang memuat dua raja. ``Sekurang-kurangnya'' hampir selalu
        menuntut komplemen
        (\cref{ex:b1:counting:complement}).
\end{enumerate}
\end{remark}

\begin{example}[Cacahan bergaya poker]\label{ex:b1:counting:poker}
Dari setumpuk $52$ kartu, banyaknya tangan berisi $5$ kartu adalah
$\binom{52}{5} = 2\,598\,960$. Tangan yang memuat tepat satu raja:
pilih rajanya ($4$ cara) lalu $4$ kartu di antara $48$ kartu bukan raja:
$4 \binom{48}{4} = 778\,320$. Kaidah hasil kali berlaku karena
pilihannya terbelah atas tahap-tahap yang saling bebas.
\end{example}

\begin{method}[Pencacahan ganda]\label{met:b1:counting:doublecount}
Untuk membuktikan kesamaan antara dua ungkapan cacahan, carilah satu
\omterm{def:b1:counting:card}{himpunan hingga} yang dicacah oleh kedua ruasnya --- lazimnya \omterm{def:b1:logic:sets}{himpunan} \emph{pasangan}
--- lalu hitung kardinalitasnya dengan dua urutan yang berbeda.
Prototipenya adalah \emph{lema jabat tangan}: pada sebuah pesta, cacah pasangan
(orang, tangan yang dijabat). Menjumlahkan atas orang memberikan $\sum_p d_p$ (banyaknya
jabat tangan tiap orang $p$); menjumlahkan atas jabat tangan memberikan
dua kali banyaknya jabat tangan (masing-masing melibatkan dua orang). Jadi
$\sum_p d_p$ genap --- sehingga banyaknya orang yang menjabat tangan
sejumlah ganjil selalu genap, kesimpulan tak sepele yang diperoleh
tanpa rumus sama sekali. Mesin yang sama menggerakkan
\cref{exo:b1:counting:12} dan beberapa pertanyaan pada soal akhir
pekan di bawah.
\end{method}

\begin{example}[Himpunan bagian rata-rata]\label{ex:b1:counting:averagesubset}
Berapa \omterm{def:b1:counting:card}{kardinalitas} rata-rata sebuah \omterm{def:b1:logic:sets}{himpunan} bagian dari \omterm{def:b1:logic:sets}{himpunan} beranggota $n$
unsur, katakanlah $E$, bila semua $2^n$ \omterm{def:b1:logic:sets}{himpunan} bagiannya sama-sama mungkin? Cacah ganda
pasangan $(A, a)$ dengan $a \in A$: menjumlahkan atas \omterm{def:b1:logic:sets}{himpunan} bagian memberikan
$\sum_A \abs A$, yaitu total yang kita cari; menjumlahkan atas unsur memberikan
$n \cdot 2^{n-1}$ (masing-masing dari $n$ unsur itu terletak tepat di separuh
\omterm{def:b1:logic:sets}{himpunan} bagiannya --- pasangkan setiap $A$ yang memuat $a$ dengan $A
\setminus \{a\}$). Jadi
\[
\frac{1}{2^n}\sum_{A \subseteq E} \abs A
= \frac{n\,2^{n-1}}{2^n} = \frac n2 :
\]
\omterm{def:b1:logic:sets}{himpunan} bagian, rata-rata, terisi separuh --- sebagaimana juga diramalkan simetri $A
\leftrightarrow \overline A$ (yang memasangkan ukuran $k$ dengan $n - k$).
Dua bukti, satu jawaban, dan keduanya menghindari perhitungan langsung
$\sum_k k\binom nk$ pada \cref{exo:b1:counting:5}: pemasangan yang
dipilih dengan baik sering menggantikan sebuah kesamaan.
\end{example}

\section{Inklusi--eksklusi}

\begin{theorem}[Inklusi--eksklusi]\label{thm:b1:counting:inclexcl}
Untuk \omterm{def:b1:counting:card}{himpunan hingga} $A_1, \dots, A_p$:
\[
\Bigl|\, \bigcup_{i=1}^{p} A_i \,\Bigr|
= \sum_{\emptyset \neq I \subseteq \intint{1}{p}}
(-1)^{\abs{I}+1} \Bigl|\, \bigcap_{i \in I} A_i \,\Bigr| .
\]
Untuk $p = 3$:
$\abs{A \cup B \cup C} = \abs A + \abs B + \abs C - \abs{A \cap B} -
\abs{A \cap C} - \abs{B \cap C} + \abs{A \cap B \cap C}$.
\end{theorem}

\begin{proof}
Tetapkan sebuah unsur $x$ pada gabungannya lalu cacah sumbangannya pada
ruas kanan. Misalkan $J = \{i : x \in A_i\}$, berkardinalitas $m \geq
1$. Unsur $x$ tercacah sekali pada $\abs{\bigcap_{i \in I} A_i}$
tepat ketika $\emptyset \neq I \subseteq J$, dengan tanda
$(-1)^{\abs I + 1}$; sumbangan totalnya adalah
\[
\sum_{k=1}^{m} \binom{m}{k} (-1)^{k+1}
= 1 - \sum_{k=0}^{m} \binom mk (-1)^k = 1 - 0 = 1
\]
menurut \cref{ex:b1:counting:binomial}. Jadi setiap unsur gabungan itu
tercacah tepat sekali.
\end{proof}

\begin{example}[Mencacah bilangan bulat yang saling prima]\label{ex:b1:counting:coprime}
Berapa banyak bilangan bulat di $\intint1{120}$ yang saling prima dengan $120 = 2^3
\times 3 \times 5$? Sebuah bilangan bulat berbagi faktor dengan $120$ tepat
ketika ia habis dibagi $2$, $3$ atau $5$, jadi cacahlah komplemen
$A_2 \cup A_3 \cup A_5$, dengan $A_d$ mengumpulkan kelipatan
$d$. Di dalam $\intint1{120}$, kelipatan $d$ berjumlah $120/d$
setiap kali $d$ membagi $120$ --- tanpa perlu fungsi lantai --- dan
$A_2 \cap A_3 = A_6$, dan seterusnya. Inklusi--eksklusi:
\[
\abs{A_2 \cup A_3 \cup A_5}
= 60 + 40 + 24 - 20 - 12 - 8 + 4 = 88 ,
\]
jadi ada $120 - 88 = 32$ bilangan bulat yang saling prima dengan $120$. Menarik
untuk mengelompokkan ulang perhitungan itu sebagai hasil kali:
\[
120 - 88 = 120\Bigl(1 - \frac12\Bigr)\Bigl(1 -
\frac13\Bigr)\Bigl(1 - \frac15\Bigr) = 120 \cdot \frac12 \cdot
\frac23 \cdot \frac45 = 32 :
\]
menjabarkan ketiga kurung itu mereproduksi persis kedelapan
suku bertanda pada inklusi--eksklusi, satu untuk setiap \omterm{def:b1:logic:sets}{himpunan} bagian $\{2, 3,
5\}$. Bentuk hasil kali ini mendefinisikan fungsi totien Euler, yang peran
aritmetikanya muncul bersama kekongruenan \cref{ch:b1:arith} dan
dikembangkan pada jilid Tahun ke-2.
\end{example}

\begin{example}[Permutasi kacau]\label{ex:b1:counting:derangement}
\emph{\omterm{def:b1:counting:objects}{Permutasi} kacau} adalah \omterm{def:b1:counting:objects}{permutasi} tanpa titik tetap. Misalkan $A_i$
\omterm{def:b1:logic:sets}{himpunan} \omterm{def:b1:counting:objects}{permutasi} $\intint{1}{n}$ yang menetapkan $i$;
maka $\abs{\bigcap_{i \in I} A_i} = (n - \abs I)!$, dan
inklusi--eksklusi mencacah \omterm{def:b1:counting:objects}{permutasi} yang mempunyai sekurang-kurangnya satu
titik tetap; \omterm{def:b1:counting:objects}{permutasi} kacau berjumlah
\[
D_n = n! \sum_{k=0}^{n} \frac{(-1)^k}{k!} .
\]
Karena $\sum (-1)^k / k! \to \eu^{-1}$ (lihat \cref{ch:b1:series}), sekitar
$37\%$ dari semua \omterm{def:b1:counting:objects}{permutasi} adalah \omterm{def:b1:counting:objects}{permutasi} kacau, berapa pun $n$ itu.
\end{example}

\begin{remark}[Di mana bab ini dipakai]\label{rem:b1:counting:whereused}
\omterm{def:b1:counting:objects}{Koefisien binomial} adalah objek bab ini yang paling banyak dipakai ulang:
ia menggerakkan teorema binomial pada \cref{ch:b1:poly}
(penjabaran $(X + a)^n$), rumus Leibniz bagi turunan ke-$n$ sebuah
hasil kali pada \cref{ch:b1:derivative}, dan
koefisien ekspansi Taylor pada \cref{ch:b1:taylor}.
\omterm{def:b1:counting:objects}{Permutasi} kembali sebagai grup --- dengan tanda yang dibangun dari
pencacahan inversi --- pada \cref{ch:b1:structures}, dan
tanda itu pada gilirannya mendefinisikan determinan pada \cref{ch:b1:det}.
Inklusi--eksklusi dan asas pencacahan adalah tulang punggung \omterm{def:b1:counting:card}{hingga}
bagi peluang diskret, yang dikembangkan pada jilid Tahun ke-2;
bilangan \omterm{def:b1:counting:objects}{permutasi} kacau pada \cref{ex:b1:counting:derangement}
dipelajari secara mendalam pada soal akhir pekan di bawah.
\end{remark}

\section{Latihan}

\begin{exercise}[$\star$]\label{exo:b1:counting:1}
Sebuah pelat nomor terdiri atas dua huruf (A--Z), lalu tiga angka, lalu
dua huruf lagi. Berapa banyak pelat yang mungkin? Berapa banyak yang tidak
mempunyai huruf berulang di antara keempat hurufnya?
\end{exercise}

\begin{exercise}[$\star$]\label{exo:b1:counting:2}
Berapa banyak anagram (penyusunan ulang huruf, bermakna atau tidak)
yang dimiliki kata \textsc{sepatu}? Dan \textsc{ananas}?
\end{exercise}

\begin{exercise}[$\star$]\label{exo:b1:counting:3}
Sebuah panitia beranggota $4$ orang dipilih dari $7$ perempuan dan $5$
laki-laki. Berapa banyak panitia: seluruhnya? yang beranggota tepat $2$
perempuan? yang beranggota sekurang-kurangnya satu laki-laki?
\end{exercise}

\begin{exercise}[$\star$]\label{exo:b1:counting:4}
Buktikan bahwa pada sebarang kelompok berisi $13$ orang, ada dua orang yang sama bulan lahirnya;
dan bahwa di antara sebarang $n + 1$ bilangan bulat yang dipilih dari $\intint{1}{2n}$, ada dua yang berurutan. \emph{(Sarang merpati dua-duanya: sebutkan
kotaknya.)}
\end{exercise}

\begin{exercise}[$\star$]\label{exo:b1:counting:5}
Hitung $\sum_{k=0}^{n} k \binom{n}{k}$. \emph{Petunjuk: turunkan
$(1 + x)^n$, atau pakai $k \binom nk = n \binom{n-1}{k-1}$ (buktikan dulu).}
\end{exercise}

\begin{exercise}[$\star\star$]\label{exo:b1:counting:6}
Ada berapa banyak \omterm{def:b1:logic:map}{pemetaan} naik tegas dari $\intint{1}{k}$ ke $\intint{1}{n}$? Simpulkan banyaknya \omterm{def:b1:logic:map}{pemetaan}
naik (tidak harus tegas). \emph{Petunjuk untuk cacahan kedua:
$f$ naik $\mapsto$ $g(i) = f(i) + i - 1$.}
\end{exercise}

\begin{exercise}[$\star\star$]\label{exo:b1:counting:7}
(Vandermonde) Buktikan, dengan mencacah \omterm{def:b1:logic:sets}{himpunan} bagian beranggota $k$ dari sebuah \omterm{def:b1:logic:sets}{himpunan} yang terbelah atas dua
blok berukuran $m$ dan $n$:
\[
\binom{m+n}{k} = \sum_{j=0}^{k} \binom{m}{j} \binom{n}{k-j} .
\]
Simpulkan $\sum_{j=0}^{n} \binom nj^2 = \binom{2n}{n}$.
\end{exercise}

\begin{exercise}[$\star\star$]\label{exo:b1:counting:8}
Berapa banyak bilangan bulat di $\intint{1}{1000}$ yang habis dibagi
$2$ atau $3$ atau $5$? (Inklusi--eksklusi; $\lfloor 1000/6 \rfloor$
mencacah kelipatan $6$, dan seterusnya.)
\end{exercise}

\begin{exercise}[$\star\star$]\label{exo:b1:counting:9}
Cacah surjeksi dari \omterm{def:b1:logic:sets}{himpunan} beranggota $4$ unsur pada \omterm{def:b1:logic:sets}{himpunan} beranggota $2$
unsur; lalu pada \omterm{def:b1:logic:sets}{himpunan} beranggota $3$ unsur. \emph{Petunjuk: cacah
\omterm{def:b1:logic:map}{pemetaan} yang tak \omterm{def:b1:logic:inj}{surjektif} dengan inklusi--eksklusi atas nilai yang terlewat.}
\end{exercise}

\begin{exercise}[$\star\star$]\label{exo:b1:counting:10}
(Bintang dan sekat) Buktikan bahwa banyaknya pemilihan-$k$ dari $n$
objek \emph{dengan} pengulangan dan urutan diabaikan --- setara dengan
banyaknya $(x_1, \dots, x_n) \in \N^n$ yang memenuhi $x_1 + \dots + x_n = k$
--- adalah $\binom{n + k - 1}{k}$. \emph{Petunjuk: sandikan sebuah penyelesaian sebagai sebaris
$k$ bintang dan $n - 1$ sekat.}
\end{exercise}

\begin{exercise}[$\star\star\star$]\label{exo:b1:counting:11}
Buktikan rumus \cref{ex:b1:counting:derangement} untuk $D_n$ secara
terperinci, lalu simpulkan $n! = \sum_{k=0}^{n} \binom{n}{k} D_{n-k}$ (buktikan pula
kesamaan ini secara langsung dengan menggolongkan \omterm{def:b1:counting:objects}{permutasi} menurut
\omterm{def:b1:logic:sets}{himpunan} titik tetapnya).
\end{exercise}

\begin{exercise}[$\star\star\star$]\label{exo:b1:counting:12}
Untuk $n \in \N^*$, buktikan dengan pencacahan ganda pasangan (\omterm{def:b1:logic:sets}{himpunan} bagian, unsur
bertanda):
\[
\sum_{k=1}^{n} k \binom{n}{k} = n\, 2^{n-1},
\qquad\text{lalu}\qquad
\sum_{k=1}^{n} k^2 \binom{n}{k} = n(n+1)\, 2^{n-2} .
\]
\emph{Untuk yang kedua: cacah pasangan unsur bertanda, sama atau tidak.}
\end{exercise}

\section{Soal: Permutasi kacau, atau surat yang salah alamat}

\begin{problem}\label{pb:b1:counting:1}
Seorang sekretaris memasukkan $n$ surat ke dalam $n$ amplop beralamat secara acak:
berapa peluang bahwa \emph{tak seorang pun} menerima surat yang benar?
Pertanyaan klasik ini (Montmort, 1708) mengantar ke bilangan \omterm{def:b1:counting:objects}{permutasi}
kacau $D_n$ pada \cref{ex:b1:counting:derangement}. Rumus
inklusi--eksklusi hanyalah langkah pembuka: soal ini
mengembangkan rekursi yang menghitung $D_n$, dua bukti mandiri lain
bagi rumus itu, teorema mencengangkan bahwa $D_n$ adalah
bilangan bulat terdekat dengan $n!/\eu$, distribusi lengkap titik tetap
sebuah \omterm{def:b1:counting:objects}{permutasi} acak, serta aritmetika ganjil barisan
$(D_n)$.\index{permutasi kacau}
Di sepanjang soal ini, $D_n$ menyatakan banyaknya \omterm{def:b1:counting:objects}{permutasi} kacau (\omterm{def:b1:counting:objects}{permutasi}
tanpa titik tetap) dari $\intint1n$, dengan kesepakatan $D_0 = 1$
(\omterm{def:b1:counting:objects}{permutasi} kosong tidak mempunyai titik tetap).

\medskip
\noindent\textbf{Bagian I --- Kasus kecil dan sensus titik
tetap.}
\begin{enumerate}
  \item Hitung $D_1, D_2, D_3$ secara langsung, dan $D_4$ dengan mendaftar
        \omterm{def:b1:counting:objects}{permutasi} kacau $\{1, 2, 3, 4\}$ yang dikelompokkan menurut nilai
        $\sigma(1)$. (Anda semestinya memperoleh $D_4 = 9$.)
  \item Untuk $0 \leq k \leq n$, tunjukkan bahwa banyaknya $P_k(n)$
        \omterm{def:b1:counting:objects}{permutasi} $\intint1n$ yang mempunyai \emph{tepat} $k$ titik
        tetap adalah $\binom nk D_{n-k}$.
  \item Periksa sensusnya untuk $n = 4$: hitung $P_0(4), \dots,
        P_4(4)$ lalu periksa bahwa jumlahnya $4! = 24$. Mana yang lebih
        mungkin untuk empat surat: tak ada yang cocok, atau tepat satu yang cocok?
  \item Dengan pencacahan ganda (\cref{met:b1:counting:doublecount})
        pasangan $(\sigma, i)$ yang memenuhi $\sigma(i) = i$, tunjukkan bahwa
        \[
        \sum_{\sigma} \abs{\mathrm{Fix}(\sigma)} = n! :
        \]
        rata-rata, sebuah \omterm{def:b1:counting:objects}{permutasi} acak mempunyai \emph{tepat satu}
        titik tetap, berapa pun $n \geq 1$ itu.
\end{enumerate}

\medskip
\noindent\textbf{Bagian II --- Dua rekursi dan dua bukti baru bagi
rumus itu.}
\begin{enumerate}[resume]
  \item Buktikan secara kombinatorial, untuk $n \geq 1$:
        \[
        D_{n+1} = n\,(D_n + D_{n-1}) .
        \]
        (Golongkan \omterm{def:b1:counting:objects}{permutasi} kacau $\sigma$ dari $\intint1{n+1}$ menurut
        $j = \sigma(n+1)$, lalu menurut apakah $\sigma(j) = n + 1$; pada
        kasus $\sigma(j) \neq n+1$, bangunlah bijeksi dengan
        \omterm{def:b1:counting:objects}{permutasi} kacau $\intint1n$ dengan mengalihkan \omterm{def:b1:logic:map}{prapeta}
        $n + 1$ ke $j$.) Periksa rekursi itu secara numerik sampai
        $D_6$.
  \item Dengan menulis $u_n = D_n - n D_{n-1}$, simpulkan dari pertanyaan 5
        bahwa $u_{n+1} = -u_n$, lalu simpulkan rekursi kedua:
        \[
        D_n = n D_{n-1} + (-1)^n \qquad (n \geq 1).
        \]
  \item Dari pertanyaan 6, buktikan dengan induksi rumus pada
        \cref{ex:b1:counting:derangement},
        \[
        D_n = n! \sum_{k=0}^{n} \frac{(-1)^k}{k!},
        \]
        --- sebuah bukti yang sama sekali tak bergantung pada inklusi--eksklusi.
  \item (Inversi binomial) Misalkan $(a_n)$ dan $(b_n)$ dua
        barisan sedemikian sehingga $a_n = \sum_{k=0}^n \binom nk b_k$ untuk
        setiap $n$. Buktikan bahwa
        \[
        b_n = \sum_{k=0}^{n} (-1)^{n-k} \binom nk a_k
        \qquad (n \in \N).
        \]
        (Tegakkan lebih dulu \emph{revisi trinomial}
        $\binom nk \binom kj = \binom nj \binom{n-j}{k-j}$, lalu
        pakai jumlah baris berselang-seling pada
        \cref{ex:b1:counting:binomial}.)
  \item Terapkan pertanyaan 8 pada kesamaan $n! = \sum_k \binom nk
        D_{n-k}$ dari \cref{exo:b1:counting:11} untuk memperoleh bukti
        \emph{ketiga} bagi rumus $D_n$.
\end{enumerate}

\medskip
\noindent\textbf{Bagian III --- Bilangan bulat terdekat dengan
$n!/\eu$.} Terima untuk bagian ini --- teorinya dibangun pada
\cref{ch:b1:series} --- bahwa
$\eu^{-1} = \lim_{n \to \infty} s_n$ dengan $s_n = \sum_{k=0}^{n}
\frac{(-1)^k}{k!}$, beserta batas deret berselang-seling yang tegas
$\abs{\eu^{-1} - s_n} < \frac1{(n+1)!}$ untuk setiap $n$.
\begin{enumerate}[resume]
  \item Tunjukkan bahwa $\bigl| D_n - n!/\eu \bigr| < \frac1{n+1}$ untuk
        setiap $n \in \N$.
  \item Simpulkan teorema utamanya: \emph{untuk setiap $n \geq 1$,
        $D_n$ adalah bilangan bulat terdekat dengan $n!/\eu$}. Mengapa
        argumen itu memerlukan $n \geq 1$?
  \item Tentukan tanda galatnya: tunjukkan bahwa $D_n > n!/\eu$
        tepat ketika $n$ genap. (Cari suku pertama yang diabaikan
        pada deret berselang-seling itu.)
  \item Hitung $D_7$ sampai $D_{10}$ dengan rekursi
        pertanyaan 5, lalu cocokkan $D_{10}$ dengan $10!/\eu$
        ($10! = 3\,628\,800$, $\eu \approx 2.718281828$).
  \item (Peluang penitipan topi) Misalkan $p_n = D_n/n!$ peluang
        bahwa sebuah \omterm{def:b1:counting:objects}{permutasi} acak seragam merupakan
        \omterm{def:b1:counting:objects}{permutasi} kacau. Tunjukkan $\abs{p_n - \eu^{-1}} < \frac1{(n+1)!}$
        lalu hitung $p_6$ sampai lima angka desimal. Beri komentar: mengapa
        jawaban atas pertanyaan Montmort pada dasarnya tidak bergantung
        pada $n$ --- bahkan sudah untuk selusin surat?
\end{enumerate}

\medskip
\noindent\textbf{Bagian IV --- Distribusi titik tetap.}
\begin{enumerate}[resume]
  \item Tetapkan $k \in \N$. Tunjukkan bahwa proporsi \omterm{def:b1:counting:objects}{permutasi}
        $\intint1n$ yang mempunyai tepat $k$ titik tetap memenuhi
        \[
        \frac{P_k(n)}{n!} = \frac{s_{n-k}}{k!}
        \;\xrightarrow[n \to \infty]{}\; \frac{\eu^{-1}}{k!} .
        \]
        (Nilai limit ini, yang berjumlah $1$, membentuk
        \emph{distribusi Poisson} berparameter $1$, objek pusat
        pada mata kuliah peluang di jilid Tahun ke-2.)
  \item Dengan mencacah ganda tripel $(\sigma, i, j)$ yang $i
        \neq j$-nya sama-sama ditetapkan oleh $\sigma$, tunjukkan bahwa
        $\sum_\sigma \abs{\mathrm{Fix}(\sigma)}\,
        (\abs{\mathrm{Fix}(\sigma)} - 1) = n!$ untuk $n \geq 2$.
        Digabungkan dengan pertanyaan 4: rata-rata
        $\abs{\mathrm{Fix}}^2$ adalah $2$, jadi ``sebaran'' (ragam)
        banyaknya titik tetap sama dengan $1$ --- sekali lagi
        tak bergantung pada $n$, dan sekali lagi cocok dengan hukum Poisson.
  \item Hitung proporsi \omterm{def:b1:counting:objects}{permutasi} yang mempunyai sekurang-kurangnya satu
        titik tetap untuk $n = 4, 5, 6$ (sebagai pecahan dan sampai empat
        angka desimal), lalu bandingkan dengan $1 - \eu^{-1} \approx
        0.6321$.
  \item Tunjukkan secara langsung --- tanpa perlu limit --- bahwa $s_{n+2} - s_n
        = (-1)^{n+1}\bigl(\frac1{(n+1)!} - \frac1{(n+2)!}\bigr)$,
        lalu simpulkan bahwa peluang $p_n = s_n$ pada
        pertanyaan 14 berayun: $p_0 > p_2 > p_4 > \dots$ dan
        $p_1 < p_3 < p_5 < \dots$, dengan nilai yang genap menurun dan
        nilai yang ganjil menaik menuju limit bersama
        $\eu^{-1}$.
  \item (Tukar kado rahasia) $n$ orang masing-masing menarik satu nama dari sebuah topi;
        jika ada yang menarik namanya sendiri, \emph{seluruh} penarikan
        diulang dari awal. Dengan memakai fakta baku bahwa sebuah
        kejadian berpeluang $p$ rata-rata menuntut $1/p$ percobaan,
        taksirlah rata-rata banyaknya penarikan lengkap yang diperlukan,
        lalu simpulkan bahwa prosedur itu menelan sekitar $\eu \approx
        2.72$ penarikan secara rata-rata, pada dasarnya tak bergantung pada
        $n$.
\end{enumerate}

\medskip
\noindent\textbf{Bagian V --- Aritmetika $D_n$, dan sebuah
sintesis.}
\begin{enumerate}[resume]
  \item Perhalus pertanyaan 5: tunjukkan bahwa untuk $j \in
        \intint2n$ yang tetap, \omterm{def:b1:counting:objects}{permutasi} kacau $\intint1n$ dengan
        $\sigma(1) = j$ berjumlah tepat $D_{n-1} + D_{n-2}$,
        tak bergantung pada $j$. Simpulkan bahwa $n - 1$ membagi $D_n$
        untuk setiap $n \geq 2$.
  \item Buktikan bahwa $D_n$ ganjil jika dan hanya jika $n$ genap. (Bekerjalah
        modulo $2$ pada rekursi pertanyaan 6.)
  \item Buktikan bahwa $D_n \equiv (-1)^n \pmod n$ untuk $n \geq 1$, lalu
        periksa kekongruenan itu pada angka terakhir $D_{10}$.
  \item Tunjukkan dari pertanyaan 6 bahwa $\dfrac{D_n}{D_{n-1}} = n +
        \dfrac{(-1)^n}{D_{n-1}}$ untuk $n \geq 3$, jadi rasio dua
        bilangan \omterm{def:b1:counting:objects}{permutasi} kacau yang berurutan \emph{hampir persis}
        $n$; jelaskan dalam satu kalimat mengapa hal ini sejalan dengan
        $D_n \approx n!/\eu$.
  \item Di mana persisnya soal ini memakai: (i) kaidah hasil kali dan
        kaidah jumlah; (ii) pencacahan ganda; (iii) teorema binomial;
        (iv) batas deret berselang-seling yang diterima tanpa bukti? Satu
        kalimat untuk masing-masing.
  \item Sintesis. Rumus untuk $D_n$ kini mempunyai tiga bukti
        (inklusi--eksklusi, rekursi beserta induksi, inversi
        binomial). Dalam satu paragraf pendek, bandingkan apa yang
        \emph{dijelaskan} masing-masing bukti: mana yang menghitung paling
        cepat, mana yang dapat diperluas ke cacahan titik tetap yang lain,
        dan mana yang mengungkap mengapa $\eu$ muncul pada soal tentang amplop.

\end{enumerate}
\end{problem}
